\section{Sujet n\textsuperscript{o}\,17}
\noindent\begin{minipage}{0.5\linewidth}\footnotesize Office du Baccalauréat du Cameroun\end{minipage}%
\hfill\begin{minipage}{0.45\linewidth}\footnotesize\raggedleft Examen : \textbf{Baccalauréat} · Série : \textbf{C/E}\\
Épreuve : \textbf{Mathématiques} · Durée : \textbf{4\,h} · Coef. : \textbf{7}\end{minipage}
\par\smallskip{\color{onbo}\rule{\linewidth}{1pt}}

\subsection*{Partie A — Évaluation des ressources \hfill (15 points)}

\noindent\textbf{Exercice 1.} \hfill\textit{(3 points)}\\
La ferme laitière « Adamaoua Lait » de Ngaoundéré possède un troupeau de $15$ vaches,
dont $6$ sont de race Goudali et $9$ d'une autre race. Pour un contrôle vétérinaire,
on choisit \emph{simultanément} $4$ vaches au hasard. Soit $X$ le nombre de vaches
Goudali parmi les $4$ choisies.
\begin{enumerate}[label=\arabic*)]
\item Justifier que le nombre total de choix possibles est $1365$. \hfill\textit{(0,5 pt)}
\item Déterminer la loi de probabilité de $X$. \hfill\textit{(1,5 pt)}
\item Calculer la probabilité d'obtenir au moins une vache Goudali. \hfill\textit{(1 pt)}
\end{enumerate}

\smallskip
\noindent\textbf{Exercice 2.} \hfill\textit{(3 points)}\\
À la laiterie, un litre de lait entier et un litre de lait écrémé n'ont pas le même
prix. Deux clients ont acheté :
\[
\left\{\begin{array}{l}
3x+2y=2700\\[2pt]
2x+5y=3450
\end{array}\right.
\]
où $x$ et $y$ désignent les prix (en F~CFA) du litre de lait entier et écrémé.
\begin{enumerate}[label=\arabic*)]
\item Écrire ce système sous la forme matricielle $M\!\begin{pmatrix}x\\y\end{pmatrix}=\begin{pmatrix}2700\\3450\end{pmatrix}$
et calculer le déterminant de $M$. \hfill\textit{(1 pt)}
\item Déterminer la matrice inverse $M^{-1}$. \hfill\textit{(1 pt)}
\item En déduire les prix $x$ et $y$. \hfill\textit{(1 pt)}
\end{enumerate}

\smallskip
\noindent\textbf{Exercice 3.} \hfill\textit{(4 points)}\\
La concentration d'un additif dans une cuve de lait (en g/L) au cours du temps $t$
(en heures) est modélisée par la fonction $f$ définie sur $[0\,;+\infty[$ par
$f(t)=(t+2)\,\mathrm{e}^{-t}$, de courbe $(\mathcal{C})$.
\begin{enumerate}[label=\arabic*)]
\item Calculer $\displaystyle\lim_{t\to+\infty}f(t)$ et interpréter graphiquement. \hfill\textit{(1 pt)}
\item Montrer que $f'(t)=-(t+1)\,\mathrm{e}^{-t}$ et dresser le tableau de variations de $f$
sur $[0\,;+\infty[$. \hfill\textit{(1,5 pt)}
\item À l'aide d'une intégration par parties, calculer
$\displaystyle J=\int_0^{1}f(t)\,\dd t$. \hfill\textit{(1,5 pt)}
\end{enumerate}

\smallskip
\noindent\textbf{Exercice 4.} \hfill\textit{(5 points)}\\
Le plan complexe est muni d'un repère orthonormé direct $(O\,;\,\vec u,\vec v)$.
On considère les points $A$ et $B$ d'affixes $z_A=2+i$ et $z_B=4+3i$.
\begin{enumerate}[label=\arabic*)]
\item Calculer le module et un argument de $z_B-z_A$. \hfill\textit{(1 pt)}
\item Soit $S$ la similitude directe d'écriture complexe $z'=(1+i)z-3i$.
Déterminer le rapport et l'angle de $S$. \hfill\textit{(1,5 pt)}
\item Déterminer l'affixe du centre $\Omega$ de la similitude $S$. \hfill\textit{(1,5 pt)}
\item Déterminer l'affixe de l'image $A'$ du point $A$ par $S$. \hfill\textit{(1 pt)}
\end{enumerate}

\subsection*{Partie B — Évaluation des compétences \hfill (5 points)}

\noindent\textit{Madame Bouba gère une ferme laitière à Ngaoundéré. Elle sollicite ton
aide, en tant qu'élève de Terminale~C, pour trois décisions. Sa production
journalière de lait, de \SI{800}{\litre} actuellement, augmente de \SI{5}{\percent}
chaque mois. Elle veut aussi construire un enclos rectangulaire \emph{adossé à un mur}
(le côté longeant le mur n'est pas clôturé) avec \SI{120}{\metre} de barrière, pour y
loger ses veaux ; elle souhaite l'aire maximale possible. Enfin, à la traite,
\SI{8}{\percent} des bidons sont mal stérilisés ; un inspecteur en contrôle $6$ au
hasard.}

\begin{enumerate}[label=\textbf{Tâche \arabic*.},leftmargin=2.4em]
\item Déterminer le mois à partir duquel la production journalière dépassera
\SI{1200}{\litre}. \hfill\textit{(1,5 pt)}
\item Déterminer les dimensions de l'enclos d'\emph{aire maximale}, ainsi que cette aire. \hfill\textit{(1,5 pt)}
\item Déterminer la probabilité qu'\emph{au moins un} des $6$ bidons contrôlés soit mal
stérilisé. \hfill\textit{(1,5 pt)}
\end{enumerate}
\par\smallskip{\footnotesize\itshape Présentation générale : 0,5 point.}

% ════════════════════════════════════════════════════
\corrige{du sujet 17}

\noindent\textbf{Partie A — Exercice 1.}
1) Choix simultané de $4$ vaches parmi $15$ : $\binom{15}{4}=\dfrac{15\cdot14\cdot13\cdot12}{24}=\mathbf{1365}$.
2) Loi hypergéométrique : $P(X{=}k)=\dfrac{\binom{6}{k}\binom{9}{4-k}}{1365}$ pour $k\in\{0,1,2,3,4\}$.
$P(0)=\frac{\binom90\binom94}{1365}=\frac{126}{1365}$ ;
$P(1)=\frac{\binom61\binom93}{1365}=\frac{6\times84}{1365}=\frac{504}{1365}$ ;
$P(2)=\frac{\binom62\binom92}{1365}=\frac{15\times36}{1365}=\frac{540}{1365}$ ;
$P(3)=\frac{\binom63\binom91}{1365}=\frac{20\times9}{1365}=\frac{180}{1365}$ ;
$P(4)=\frac{\binom64\binom90}{1365}=\frac{15}{1365}$.
(Vérification : $126+504+540+180+15=1365$.)
3) $P(X\ge1)=1-P(0)=1-\frac{126}{1365}=\frac{1239}{1365}=\dfrac{59}{65}\approx0{,}908$.

\noindent\textbf{Exercice 2.}
1) $M=\begin{pmatrix}3&2\\2&5\end{pmatrix}$, $\det M=3\times5-2\times2=15-4=\mathbf{11}$.
2) $M^{-1}=\dfrac{1}{11}\begin{pmatrix}5&-2\\-2&3\end{pmatrix}$
(formule $\frac{1}{\det M}\begin{pmatrix}d&-b\\-c&a\end{pmatrix}$).
3) $\begin{pmatrix}x\\y\end{pmatrix}=M^{-1}\begin{pmatrix}2700\\3450\end{pmatrix}
=\frac{1}{11}\begin{pmatrix}5\cdot2700-2\cdot3450\\-2\cdot2700+3\cdot3450\end{pmatrix}
=\frac{1}{11}\begin{pmatrix}13500-6900\\-5400+10350\end{pmatrix}
=\frac{1}{11}\begin{pmatrix}6600\\4950\end{pmatrix}=\begin{pmatrix}600\\450\end{pmatrix}$.
Donc $x=\mathbf{600}$ F (litre entier) et $y=\mathbf{450}$ F (litre écrémé).
\emph{Contrôle :} $3\cdot600+2\cdot450=1800+900=2700$ ;
$2\cdot600+5\cdot450=1200+2250=3450$. \checkmark

\noindent\textbf{Exercice 3.}
1) $\lim_{t\to+\infty}(t+2)\mathrm e^{-t}=0$ (croissances comparées) : l'axe des abscisses
$y=0$ est asymptote horizontale en $+\infty$.
2) $f'(t)=\mathrm e^{-t}-(t+2)\mathrm e^{-t}=(1-t-2)\mathrm e^{-t}=-(t+1)\mathrm e^{-t}$.
Sur $[0\,;+\infty[$, $t+1>0$ donc $f'(t)<0$ : $f$ est \emph{strictement décroissante},
de $f(0)=2$ vers $0$.
3) IPP ($u=t+2,\ v'=\mathrm e^{-t}$, donc $u'=1,\ v=-\mathrm e^{-t}$) :
$J=\big[-(t+2)\mathrm e^{-t}\big]_0^1+\int_0^1\mathrm e^{-t}\dd t
=\big[-(t+2)\mathrm e^{-t}\big]_0^1+\big[-\mathrm e^{-t}\big]_0^1$.
Or $\big[-(t+2)\mathrm e^{-t}\big]_0^1=-3\mathrm e^{-1}+2$ et
$\big[-\mathrm e^{-t}\big]_0^1=-\mathrm e^{-1}+1$. D'où
$J=2-3\mathrm e^{-1}+1-\mathrm e^{-1}=3-4\mathrm e^{-1}=\mathbf{3-\dfrac{4}{\mathrm e}}\approx1{,}53$.

\noindent\textbf{Exercice 4.}
1) $z_B-z_A=(4+3i)-(2+i)=2+2i$. Module $|2+2i|=2\sqrt2$ ;
argument $\arg(2+2i)=\frac{\pi}{4}$.
2) $z'=(1+i)z-3i$ : le coefficient de $z$ est $a=1+i$. Rapport
$k=|a|=\sqrt2$ ; angle $\theta=\arg(a)=\frac{\pi}{4}$.
3) Centre $\Omega$ : point fixe $\omega=(1+i)\omega-3i$, soit
$\omega-(1+i)\omega=-3i$, $\omega(1-1-i)=-3i$, $-i\,\omega=-3i$, donc
$\omega=\dfrac{-3i}{-i}=\mathbf{3}$. Le centre a pour affixe $z_\Omega=3$.
4) $A'$ : $z_{A'}=(1+i)(2+i)-3i=(2+i+2i+i^2)-3i=(2+3i-1)-3i=1+3i-3i=\mathbf{1}$.

\smallskip
\noindent\textbf{Partie B — Situation.}
\emph{Tâche 1.} Production du mois $n$ : $u_n=800\times1{,}05^{\,n-1}$. On résout
$u_n>1200$, soit $1{,}05^{\,n-1}>\frac{3}{2}$, d'où
$n-1>\frac{\ln(1{,}5)}{\ln1{,}05}\approx\frac{0{,}4055}{0{,}0488}\approx8{,}31$ : à partir du
\textbf{10\textsuperscript{e} mois} ($u_{10}=800\times1{,}05^9\approx\mathbf{1241}$ litres).

\emph{Tâche 2.} Largeur $x$ (deux côtés perpendiculaires au mur), longueur
$y=120-2x$. Aire $A(x)=x(120-2x)=120x-2x^2$ ; $A'(x)=120-4x=0$ en $x=30$. Dimensions :
$x=\SI{30}{\metre}$, $y=\SI{60}{\metre}$ ; aire maximale $30\times60=\SI{1800}{\metre\squared}$.

\emph{Tâche 3.} Loi binomiale $\mathcal{B}(6\,;0{,}08)$.
$P(\text{au moins un})=1-(0{,}92)^6\approx1-0{,}606=\mathbf{0{,}394}$, soit environ
\SI{39,4}{\percent}.
